\documentclass[12pt]{ximera}
\title{Homework 8: Lone-Chooser, Sealed Bids, and Markers}
\begin{document}
\begin{abstract}
Sections 3.4--3.6: the lone-chooser method applied to dividing a day of lab time, the
method of sealed bids for indivisible household goods, and the method of markers for a
sixty-item beanie baby collection.
\end{abstract}
\maketitle

\section*{Sections 3.4--3.6: Lone-Chooser, Sealed Bids, and Markers}

\begin{problem}
Three research groups (A, B, and C) share a lab and need to divide lab time per day.
They decide to create a schedule by fairly dividing $24$ hours among them.

Define \emph{Daytime} hours to be 8am--8pm, and \emph{Nighttime} hours to be
8pm--8am. Each is a $12$-hour block.
\begin{itemize}
\item Group A values Daytime hours as $80\%$ of the total value, Nighttime as $20\%$.
\item Group B values all time slots equally.
\item Group C values Daytime hours as $40\%$ of the total, Nighttime as $60\%$.
\end{itemize}

Group B is player 1, so they split the day in half:
\[ s_1 = \{\text{Noon--Midnight}\}, \qquad s_2 = \{\text{Midnight--Noon}\}. \]

\begin{prompt}
\textbf{(a)} Describe $s_1$ and $s_2$ each as a portion of Daytime plus a portion of
Nighttime. Enter fractions.

$s_1 = \answer{2/3}$ of Daytime $+$ $\answer{1/3}$ of Nighttime

\smallskip
$s_2 = \answer{1/3}$ of Daytime $+$ $\answer{2/3}$ of Nighttime

\begin{hint}
$s_1$ runs Noon to Midnight. How many of those hours fall in the 8am--8pm Daytime
block, out of the $12$ Daytime hours total?
\end{hint}

\begin{solution}
$s_1$ is Noon--Midnight. Of that, Noon--8pm is $8$ hours of Daytime, which is
$\frac{8}{12}=\frac23$ of the Daytime block, and 8pm--Midnight is $4$ hours of
Nighttime, which is $\frac{4}{12}=\frac13$ of the Nighttime block.

$s_2$ is Midnight--Noon: Midnight--8am is $8$ hours of Nighttime ($\frac23$ of it) and
8am--Noon is $4$ hours of Daytime ($\frac13$ of it).
\end{solution}
\end{prompt}

\begin{prompt}
\textbf{(b)} Calculate the value research group A places on $s_1$ and $s_2$, and
indicate which they will choose.

$s_1$: $\answer{60}\%$ \qquad $s_2$: $\answer{40}\%$

\smallskip
Group A chooses:
\begin{multipleChoice}
\choice[correct]{$s_1$ (Noon--Midnight)}
\choice{$s_2$ (Midnight--Noon)}
\end{multipleChoice}

\begin{solution}
Using A's values of $80\%$ for the whole Daytime block and $20\%$ for the whole
Nighttime block:
\[
s_1 = \tfrac23(80\%) + \tfrac13(20\%) = 53.\overline{3}\% + 6.\overline{6}\% = 60\%,
\]
\[
s_2 = \tfrac13(80\%) + \tfrac23(20\%) = 26.\overline{6}\% + 13.\overline{3}\% = 40\%.
\]
A takes $s_1$ at $60\%$ --- well above the $33.3\%$ that would be a fair share.
\end{solution}
\end{prompt}

\begin{prompt}
\textbf{(c)} Both group B and group A now divide their shares into three pieces of
equal value \emph{to them}, using two cuts each.

Group B divides $s_2$ (Midnight--Noon) at $\answer{4}$am and $\answer{8}$am.

\smallskip
Group A divides $s_1$ (Noon--Midnight) at $\answer{3}$pm and $\answer{6}$pm.

\begin{hint}
Group B has no time preference, so their cuts are equally spaced. Group A values their
$60\%$ share, so each piece must be worth $20\%$ to them --- and Daytime hours are worth
far more per hour to A than Nighttime hours.
\end{hint}

\begin{solution}
\textbf{Group B} values all hours equally, so they cut their $12$-hour share into
three $4$-hour pieces: Midnight--4am, 4am--8am, and 8am--Noon.

\textbf{Group A} must split their $60\%$ share into three pieces worth $20\%$ each.
For A, the Daytime block is worth $80\%$ over $12$ hours, or
$6.\overline{6}\%$ per hour, while Nighttime is worth $20\%$ over $12$ hours, or
$1.\overline{6}\%$ per hour.

Starting at Noon in Daytime, reaching $20\%$ takes
\[ 20 \div 6.\overline{6} = 3 \text{ hours}, \]
so the first cut is at 3pm and the second at 6pm. The pieces are:
\begin{itemize}
\item Noon--3pm: $3$ Daytime hours $= 20\%$;
\item 3pm--6pm: $3$ Daytime hours $= 20\%$;
\item 6pm--Midnight: $2$ Daytime hours ($13.\overline{3}\%$) plus $4$ Nighttime hours
      ($6.\overline{6}\%$) $= 20\%$.
\end{itemize}
Note how uneven the pieces are in clock time --- $3$, $3$, and $6$ hours --- precisely
because A does not value all hours equally.
\end{solution}
\end{prompt}

\begin{prompt}
\textbf{(d)} Group C now selects one piece from each of B's and A's divisions. What is
the finalized lab schedule?

First, C's valuations. From B's pieces:
Midnight--4am $= \answer{20}\%$, \quad 4am--8am $= \answer{20}\%$, \quad
8am--Noon $= \answer[tolerance=0.05]{13.3}\%$.

\smallskip
From A's pieces: Noon--3pm $= \answer{10}\%$, \quad 3pm--6pm $= \answer{10}\%$, \quad
6pm--Midnight $= \answer[tolerance=0.05]{26.7}\%$.

\smallskip
So group C takes 6pm--Midnight from A, and Midnight--4am from B, for a total of
$\answer[tolerance=0.05]{46.7}\%$.

\begin{solution}
For group C, Daytime is worth $40\%$ over $12$ hours ($3.\overline{3}\%$ per hour) and
Nighttime $60\%$ over $12$ hours ($5\%$ per hour).

From B's pieces: Midnight--4am and 4am--8am are each $4$ Nighttime hours $= 20\%$,
while 8am--Noon is $4$ Daytime hours $= 13.\overline{3}\%$. C takes one of the
Nighttime pieces --- say Midnight--4am (the other is equally good).

From A's pieces: Noon--3pm and 3pm--6pm are each $3$ Daytime hours $= 10\%$, while
6pm--Midnight is $2$ Daytime hours ($6.\overline{6}\%$) plus $4$ Nighttime hours
($20\%$) $= 26.\overline{6}\%$. C takes 6pm--Midnight.

C's total is $20\% + 26.\overline{6}\% = 46.\overline{6}\%$.

\textbf{The finalized lab schedule:}
\begin{itemize}
\item \textbf{Group A}: Noon--6pm (worth $40\%$ to A)
\item \textbf{Group B}: 4am--Noon (worth $33.\overline{3}\%$ to B)
\item \textbf{Group C}: 6pm--4am (worth $46.\overline{6}\%$ to C)
\end{itemize}
The three blocks cover all $24$ hours with no overlap, and every group receives more
than the $33.3\%$ they need --- each by their own very different standard.
\end{solution}
\end{prompt}
\end{problem}

\begin{problem}
You and a few roommates have shared an apartment for the last two years, but now it's
time to move out, and you've agreed to divvy your shared belongings using the method of
sealed bids. The values you each assign are:

\begin{center}
\begin{tabular}{|r|c|c|c|c|}\hline
 & \textbf{You} & \textbf{Sam} & \textbf{Riley} & \textbf{Jordan} \\\hline
Sofa & \$205 & \$260 & \$185 & \$215 \\
Smart TV & \$210 & \$150 & \$180 & \$175 \\
Microwave & \$55 & \$75 & \$65 & \$50 \\
Coffee Maker & \$60 & \$25 & \$90 & \$75 \\
Blender & \$60 & \$30 & \$75 & \$100 \\
Vacuum Cleaner & \$120 & \$150 & \$130 & \$105 \\
Dining Table & \$270 & \$230 & \$255 & \$300 \\\hline
\end{tabular}
\end{center}

\begin{prompt}
\textbf{(a)} Which roommate gets to keep each item? (Each item goes to the highest
bidder.)

Sofa: $\answer{Sam}$ \quad Smart TV: $\answer{You}$ \quad Microwave: $\answer{Sam}$

\smallskip
Coffee Maker: $\answer{Riley}$ \quad Blender: $\answer{Jordan}$ \quad
Vacuum Cleaner: $\answer{Sam}$ \quad Dining Table: $\answer{Jordan}$

\begin{solution}
Reading down each row for the largest bid:
\[
\begin{array}{lll}
\text{Sofa} & \text{Sam} & \$260\\
\text{Smart TV} & \text{You} & \$210\\
\text{Microwave} & \text{Sam} & \$75\\
\text{Coffee Maker} & \text{Riley} & \$90\\
\text{Blender} & \text{Jordan} & \$100\\
\text{Vacuum Cleaner} & \text{Sam} & \$150\\
\text{Dining Table} & \text{Jordan} & \$300
\end{array}
\]
\end{solution}
\end{prompt}

\begin{prompt}
\textbf{(b)} What is each roommate's ``fair share'' in this system?

You: $\$\answer{245}$ \quad Sam: $\$\answer{230}$ \quad Riley: $\$\answer{245}$ \quad
Jordan: $\$\answer{255}$

\begin{solution}
Each roommate's total valuation, divided by four:
\[
\begin{array}{lll}
\text{You}: & 980 & \to 980/4 = \$245\\
\text{Sam}: & 920 & \to 920/4 = \$230\\
\text{Riley}: & 980 & \to 980/4 = \$245\\
\text{Jordan}: & 1020 & \to 1020/4 = \$255
\end{array}
\]
\end{solution}
\end{prompt}

\begin{prompt}
\textbf{(c)} Compare each roommate's fair share with their perceived value of the items
they've won. What does each owe, or what are they owed, in the initial settlement?
Enter a positive number for money received and a negative number for money paid in.

You: $\answer{35}$ \quad Sam: $\answer{-255}$ \quad Riley: $\answer{155}$ \quad
Jordan: $\answer{-145}$

\begin{solution}
\[
\begin{array}{lcccl}
 & \text{items won} & \text{fair share} & \text{difference} & \\
\text{You} & \$210 & \$245 & -\$35 & \text{receives } \$35\\
\text{Sam} & \$485 & \$230 & +\$255 & \text{pays } \$255\\
\text{Riley} & \$90 & \$245 & -\$155 & \text{receives } \$155\\
\text{Jordan} & \$400 & \$255 & +\$145 & \text{pays } \$145
\end{array}
\]
(Sam's $\$485$ is the Sofa, Microwave, and Vacuum: $260+75+150$. Jordan's $\$400$ is
the Blender and Dining Table: $100+300$.)
\end{solution}
\end{prompt}

\begin{prompt}
\textbf{(d)} Once every roommate has items and/or compensation worth exactly their fair
share, how much money is left in the system --- that is, what is the surplus?

$\$\answer{210}$ \qquad and each roommate's quarter of it is $\$\answer{52.50}$

\begin{solution}
Money in from Sam and Jordan is $255 + 145 = \$400$; money out to You and Riley is
$35 + 155 = \$190$. The surplus is
\[ \$400 - \$190 = \$210, \]
which splits four ways as $\$210/4 = \$52.50$ each.
\end{solution}
\end{prompt}

\begin{prompt}
\textbf{(e)} After splitting the surplus equally, what is each roommate's cash position
in the final settlement? Enter a positive number for money received and a negative
number for money paid.

You: $\answer{87.50}$ \quad Sam: $\answer{-202.50}$ \quad Riley: $\answer{207.50}$ \quad
Jordan: $\answer{-92.50}$

\begin{solution}
Add each roommate's \$52.50 share of the surplus to their initial settlement:
\[
\begin{array}{ll}
\text{You}: & +35 + 52.50 = \text{receives } \$87.50 \text{, plus the Smart TV}\\
\text{Sam}: & -255 + 52.50 = \text{pays } \$202.50 \text{, keeps Sofa, Microwave, Vacuum}\\
\text{Riley}: & +155 + 52.50 = \text{receives } \$207.50 \text{, plus the Coffee Maker}\\
\text{Jordan}: & -145 + 52.50 = \text{pays } \$92.50 \text{, keeps Blender and Dining Table}
\end{array}
\]
The cash balances: $-87.50 + 202.50 - 207.50 + 92.50 = 0$. \checkmark

Every roommate ends up with more than their fair share --- each is ahead by exactly
\$52.50 in their own valuation. That surplus exists because the roommates disagree
about what things are worth.
\end{solution}
\end{prompt}
\end{problem}

\begin{problem}
Six grandchildren are dividing up their grandma's Ty beanie baby collection using the
method of markers. They lay the $60$ beanie babies out in a randomly ordered array,
numbered $1$ through $60$, then each individually marks off six consecutive segments
that are approximately equally valuable to them.

\begin{center}
\begin{tabular}{r|c|c|c|c|c|c}
& 1st share & 2nd share & 3rd share & 4th share & 5th share & 6th share \\\hline
Allan  & 1--7  & 8--14  & 15--29 & 30--36 & 37--52 & 53--60 \\
Brenna & 1--10 & 11--20 & 21--36 & 37--43 & 44--48 & 49--60 \\
Caleb  & 1--5  & 6--20  & 21--31 & 32--40 & 41--51 & 52--60 \\
Devon  & 1--8  & 9--17  & 18--25 & 26--37 & 38--53 & 54--60 \\
Emma   & 1--9  & 10--18 & 19--32 & 33--46 & 47--50 & 51--60 \\
Frank  & 1--14 & 15--21 & 22--28 & 29--40 & 41--56 & 57--60
\end{tabular}
\end{center}

Each grandchild's five markers sit at the end of their first five shares.

\begin{prompt}
\textbf{(a)} Which of their shares is each grandchild awarded via this method?

Round 1 (first markers) --- awarded to $\answer{Caleb}$, who takes items 1--$\answer{5}$.

\smallskip
Round 2 (second markers) --- $\answer{Allan}$ takes items 6--$\answer{14}$.

\smallskip
Round 3 (third markers) --- $\answer{Devon}$ takes items 15--$\answer{25}$.

\smallskip
Round 4 (fourth markers) --- $\answer{Frank}$ takes items 26--$\answer{40}$.

\smallskip
Round 5 (fifth markers) --- $\answer{Brenna}$ takes items 41--$\answer{48}$.

\smallskip
The last remaining grandchild, $\answer{Emma}$, takes items $\answer{51}$--60.

\begin{hint}
In each round, look only at the relevant marker (first markers in round 1, second
markers in round 2, and so on) among the players who have not yet been served, and
find the leftmost one that lies past the items already handed out.
\end{hint}

\begin{solution}
The markers sit at the ends of the first five shares:
\[
\begin{array}{ll}
\text{Allan}: & 7,\ 14,\ 29,\ 36,\ 52\\
\text{Brenna}: & 10,\ 20,\ 36,\ 43,\ 48\\
\text{Caleb}: & 5,\ 20,\ 31,\ 40,\ 51\\
\text{Devon}: & 8,\ 17,\ 25,\ 37,\ 53\\
\text{Emma}: & 9,\ 18,\ 32,\ 46,\ 50\\
\text{Frank}: & 14,\ 21,\ 28,\ 40,\ 56
\end{array}
\]

\textbf{Round 1:} first markers are $7, 10, 5, 8, 9, 14$. Caleb's is leftmost at $5$,
so Caleb takes items $1$--$5$ and leaves.

\textbf{Round 2:} among the rest, second markers are Allan $14$, Brenna $20$,
Devon $17$, Emma $18$, Frank $21$. Allan's is leftmost, so Allan takes $6$--$14$.

\textbf{Round 3:} third markers are Brenna $36$, Devon $25$, Emma $32$, Frank $28$.
Devon takes $15$--$25$.

\textbf{Round 4:} fourth markers are Brenna $43$, Emma $46$, Frank $40$. Frank takes
$26$--$40$.

\textbf{Round 5:} fifth markers are Brenna $48$ and Emma $50$. Brenna takes $41$--$48$.

\textbf{Emma}, the last one left, takes everything past her final marker at $50$:
items $51$--$60$.
\end{solution}
\end{prompt}

\begin{prompt}
\textbf{(b)} Which beanie babies remain after every grandchild has been awarded exactly
their fair share?

Items $\answer{49}$ and $\answer{50}$ are left over.

\begin{solution}
The awarded segments are $1$--$5$, $6$--$14$, $15$--$25$, $26$--$40$, $41$--$48$, and
$51$--$60$. The only numbers not covered are $49$ and $50$.
\end{solution}
\end{prompt}

\begin{prompt}
\textbf{(c)} Describe one way that the rest of the beanie babies can be fairly
distributed.

\begin{freeResponse}
\end{freeResponse}

\begin{solution}
Every grandchild already holds a share they consider fair, so anything done with the
two leftovers only improves someone's position --- there is no way to make the division
unfair at this stage. Reasonable options include:

\begin{itemize}
\item \textbf{Run the method again} on the leftover items. With only two items and six
      players this is awkward, but it is the textbook answer when many items remain.
\item \textbf{Draw lots}, giving each leftover to a randomly chosen grandchild, with no
      one receiving two.
\item \textbf{Auction them} among the six, splitting the proceeds equally --- effectively
      the method of sealed bids applied to the remainder.
\item \textbf{Let the grandchildren take turns choosing}, in an order decided randomly.
\end{itemize}

Any of these is acceptable. What matters is that the leftovers are distributed by some
process the players agree to in advance, since the fairness guarantee of the method of
markers has already been satisfied.
\end{solution}
\end{prompt}
\end{problem}

\end{document}
